\section{Tìm kiếm theo chiều sâu (Depth-First Search, DFS)}
\label{sec:bloxorz-dfs}

\subsection{Biến thể đã chọn: DFS lặp với ngăn xếp}
DFS được cài đặt dưới dạng lặp, dùng ngăn xếp LIFO tường minh thay cho đệ quy.
Lớp \texttt{StackFrontier} kế thừa \texttt{Frontier} của vòng lặp chung và chỉ
đổi thao tác lấy nút: nút được đưa vào gần nhất được mở rộng trước. Cách này
tránh giới hạn độ sâu đệ quy của Python trên màn lớn và giữ nguyên chính sách
trạng thái lặp, bộ đếm cùng các giới hạn của hạ tầng dùng chung. Bộ giải vẫn
theo giao diện \texttt{solve(stage, caps)} đã đóng băng.

Nhóm chọn DFS thay vì IDS và ghi lại lựa chọn này. Trên mười màn chính thức,
DFS trả tổng $701$ bước, còn phép duyệt FIFO trong kiểm thử trả $328$ bước,
tức $2{,}1$ lần; mức tệ nhất là L4-02 với $168$ bước so với $24$. Độ dài tăng
đều như vậy vẫn đọc được như một đặc điểm của chiến lược duyệt sâu, không phải
hiện tượng đường đi dài bất thường, nên việc mở rộng lại của IDS không cần
thiết để so sánh ở mục thực nghiệm.

\subsection{Trạng thái lặp, tính kết thúc và tính tối ưu}
Mỗi trạng thái được đưa vào ngăn xếp nhiều nhất một lần. \texttt{Frontier.push}
từ chối khóa đã gặp, kể cả khi nút tương ứng còn đang chờ hoặc đã được mở rộng,
nên tập đã gặp chính là cơ chế chặn chu trình. Trên đồ thị hữu hạn, ngăn xếp
cạn dần và tìm kiếm kết thúc bằng nghiệm hoặc kết luận không có nghiệm; không
nhánh nào quay lại vô hạn.

DFS đầy đủ nhưng không tối ưu. Nó trả nghiệm đầu tiên chạm đích theo thứ tự
ngăn xếp, nên cả số bước lẫn chi phí đều có thể lớn hơn nghiệm tối ưu. Trên
L1-02, DFS trả $50$ bước với chi phí $50$, trong khi UCS ở
mục~\ref{sec:bloxorz-ucs} trả $22$ bước với chi phí $22$. Vì vậy DFS chỉ dùng
để kiểm chứng yêu cầu về duyệt sâu, không dùng thay UCS hay A* khi cần nghiệm
tối ưu.

\subsection{Kiểm chứng trên bộ dữ liệu}
Kiểm thử phát lại mọi đường đi qua \texttt{Board} và so chi phí với chi phí
tính lại, nên đường không hợp lệ hoặc chi phí sai đều bị bắt. Ba màn nhỏ có
đáp án tính tay đều được giải với $6$, $4$ và $4$ bước, khớp đáp án. Vì DFS
không bảo đảm tối ưu, kiểm thử chỉ ràng buộc độ dài không nhỏ hơn đáp án tối
ưu chứ không khẳng định bằng.

Hai bàn tổng hợp kiểm tra đúng phần đặc trưng của DFS. Bàn thứ nhất có chu
trình nhưng không có đích: tìm kiếm trả \texttt{unsolvable} sau khi mở rộng
đúng ba trạng thái, thay vì quay vòng. Bàn thứ hai có hai đường tới đích: DFS
đi đường $5$ bước trong khi vòng lặp FIFO đi đường $2$ bước, xác nhận thứ tự
LIFO chứ không chỉ xác nhận có tìm ra nghiệm.

Trên mười màn chính thức, DFS giải được cả mười và mọi đường đi phát lại hợp
lệ. DFS mở rộng ít nút hơn UCS ở $8/10$ màn vì dừng ngay khi chạm đích, nhưng
đỉnh ngăn xếp lại lớn hơn ở cả $10/10$ màn.

\subsection{Giới hạn và phạm vi kết luận}
Các số trên là kết quả của bộ giải, không phải phép đo thời gian hay bộ nhớ;
phép đo thuộc công cụ chung và chưa được thực hiện ở đây.

Việc ngăn xếp đạt đỉnh lớn hơn không đủ để kết luận DFS dùng nhiều bộ nhớ hơn:
cả DFS lẫn UCS đều giữ tập trạng thái đã gặp với $O(V)$ bộ nhớ, và đỉnh frontier
chỉ là một phần của mức dùng thực tế. Lý thuyết DFS chỉ giữ một nhánh không còn
đúng khi bắt buộc phải có tập đã gặp. Với $V$ trạng thái và $E$ cạnh đạt được,
thời gian là $O(V+E)$ vì mỗi trạng thái được mở rộng nhiều nhất một lần; bộ nhớ
là $O(V)$ cho tập đã gặp cộng với ngăn xếp.
